\documentclass[11pt,a4paper]{article}

\usepackage[utf8]{inputenc}
\usepackage[T1]{fontenc}
% babel nur laden, wenn das Sprachpaket vorhanden ist, damit das Dokument
% auch auf einer minimalen TeX-Live uebersetzt.
\IfFileExists{ngerman.ldf}{\usepackage[ngerman]{babel}}{%
  \typeout{^^JHINWEIS: babel-german fehlt, Silbentrennung bleibt englisch.^^J}}
\usepackage{amsmath,amssymb}
\usepackage{booktabs}
\usepackage{longtable}
\usepackage[margin=2.5cm]{geometry}
\usepackage{hyperref}
\usepackage{xcolor}   % blue!50!black in hypersetup braucht das Mischmodell
\usepackage{enumitem}

\hypersetup{
  colorlinks=true,
  linkcolor=blue!50!black,
  urlcolor=blue!50!black,
  pdftitle={FEMBEM -- Harmonische FEM, Stand und Befunde},
}

\newcommand{\Ar}{\mathbf{A}_r}
\newcommand{\Ai}{\mathbf{A}_i}
\newcommand{\Js}{\mathbf{J}_s}

\title{Harmonische FEM für TEAM 28\\[0.3em]
       \large Formulierung, Löser und gemessene Befunde}
\author{FEMBEM}
\date{September 2026}

\begin{document}
\maketitle

\begin{abstract}
\noindent
Dieser Bericht fasst den Stand des Pfades \texttt{--fem-direct} zusammen: eine
harmonische Wirbelstromrechnung in reiner FEM mit mitvernetztem Luftgebiet und
Sparse-Direktlöser, ohne BEM-Kopplung. Beschrieben werden die Formulierung, die
Kantenelemente, die drei Projektionen und Regularisierungen, die beiden Löser
und die Kraftauswertung. Der zweite Teil nennt die Messergebnisse mit ihren
Vorbehalten. Alle Zahlen im Text sind Läufe, keine Schätzungen; wo etwas
geschätzt ist, steht es dabei. Der Kernbefund: die Levitationskraft in der
Gleichgewichtslage beträgt 1,02\,N gegen ein Plattengewicht von
1,026\,N, aber die Netzreihe belegt diese Übereinstimmung noch nicht.
\end{abstract}

\tableofcontents

\section{Aufgabe und Abgrenzung}

TEAM~28 ist eine axialsymmetrische Levitationsanordnung: zwei konzentrische
Ringspulen mit gegenläufigem 50\,Hz-Strom, darüber eine
Aluminiumscheibe (3\,mm dick, Außenradius 65\,mm,
$\sigma = 3.4e7\,S/m$). Gesucht ist die zeitgemittelte
Auftriebskraft bei ruhender Platte, aus der sich die Schwebehöhe ergibt.

Die Aufgabe ist linear und harmonisch. Der hier beschriebene Pfad nutzt das
aus und rechnet direkt im Frequenzbereich, statt einen Einschwingvorgang zu
integrieren. Er verzichtet außerdem auf die BEM-Kopplung des Hauptpfades: das
Luftgebiet wird mitvernetzt und am Rand mit $\mathbf{A}\times\mathbf{n} = 0$
abgeschnitten. Das kostet Freiheitsgrade, erspart aber die Randintegrale und
macht das System sparse und direkt lösbar.

Zwei Einschränkungen gelten durchweg. Erstens $\mu_r = 1$ überall -- weder
Platte noch Spulen sind magnetisch. Zweitens ruht die Platte; die Kraft wird
für eine feste Lage berechnet, nicht die Bewegung integriert.

\section{Formulierung}

\subsection{Kontinuierliches Problem}

Im magnetischen Vektorpotential $\mathbf{A}$ mit $\mathbf{B} =
\nabla\times\mathbf{A}$ lautet die Wirbelstromgleichung im Frequenzbereich
\begin{equation}
  \nabla\times\left(\nu\,\nabla\times\hat{\mathbf{A}}\right)
  + \mathrm{j}\omega\sigma\hat{\mathbf{A}} = \hat{\Js},
  \qquad \nu = \frac{1}{\mu_0},
  \label{eq:kont}
\end{equation}
mit dem Phasor $\hat{\mathbf{A}}$ und der eingeprägten Spulenstromdichte
$\hat{\Js}$. Die induzierte Stromdichte ist
$\hat{\mathbf{J}} = -\mathrm{j}\omega\sigma\hat{\mathbf{A}}$.

Gleichung~\eqref{eq:kont} fixiert die Eichung nicht. In leitenden Gebieten
($\sigma > 0$) erledigt der Term $\mathrm{j}\omega\sigma\hat{\mathbf{A}}$ das
implizit, in nichtleitenden nicht: dort ist jedes Gradientenfeld
$\nabla\varphi$ eine Lösung des homogenen Problems. Abschnitt~\ref{sec:regul}
behandelt das.

\subsection{Diskretisierung und reelle Blockform}

Nach Galerkin-Diskretisierung mit Kantenelementen (Abschnitt~\ref{sec:elem})
entstehen zwei Matrizen,
\begin{equation}
  K_{ij} = \int_\Omega \nu\,(\nabla\times\mathbf{W}_i)\cdot(\nabla\times\mathbf{W}_j)\,\mathrm{d}V,
  \qquad
  C_{ij} = \int_\Omega \sigma\,\mathbf{W}_i\cdot\mathbf{W}_j\,\mathrm{d}V,
\end{equation}
und das komplexe System $(K + \mathrm{j}\omega C)\,\hat{x} = \hat{b}$. Gelöst
wird nicht komplex, sondern in der reellen $2n$-Blockform
\begin{equation}
  \underbrace{\begin{bmatrix} K & -\omega C \\ \omega C & K \end{bmatrix}}_{H}
  \begin{bmatrix} x_r \\ x_i \end{bmatrix}
  = \begin{bmatrix} b_r \\ b_i \end{bmatrix}.
  \label{eq:block}
\end{equation}
Das erspart komplexe Kernel, komplexe Löserpfade und eine zweite
Vorkonditionierung. $H$ ist \emph{nicht} symmetrisch.

$K$ und $C$ werden nicht getrennt assembliert, sondern aus zwei Aufrufen
derselben Assemblierung gewonnen: der implizite Euler-Schritt liefert
$K + C/\Delta t$, also
\begin{equation}
  \Delta t = 1/\omega \;\Rightarrow\; K + \omega C,
  \qquad
  \Delta t \to \infty \;\Rightarrow\; K,
  \qquad
  C = \frac{(K+\omega C) - K}{\omega}.
\end{equation}
Der Umweg vermeidet eine zweite vollständige Konfiguration, die von der ersten
abweichen könnte, ohne dass es auffällt.

\subsection{Anregung und Phasenkonvention}

Die Spulen sind stromgetrieben mit $i(t) = I\sin(\omega t)$. Für die
Blockform~\eqref{eq:block} wird die \emph{Amplitude} gebraucht, nicht eine
Zeitprobe: der Quellterm wird deshalb bei $t = T/4$ ausgewertet, wo
$\sin(\omega t) = 1$. Das entspricht der Konvention $i(t) = I\cos(\omega t)$
mit reeller rechter Seite, also $b_i = 0$. Für die zeitgemittelte Kraft ist
das gleichgültig -- sie ist unter einer gemeinsamen Phasendrehung aller
Phasoren invariant.

Kontrolle der Skalierung über die Amperewindungen, gerechnet aus den
protokollierten Stromdichten und der Netzgeometrie:

\begin{center}
\begin{tabular}{lrrr}
\toprule
& $J$ [A/m$^2$] & Querschnitt [m$^2$] & Amperewindungen \\
\midrule
Innenspule  & $1{,}32\cdot 10^{7}$  & $1{,}456\cdot 10^{-3}$ & $1{,}92\cdot 10^{4}$ \\
Außenspule  & $-1{,}48\cdot 10^{7}$ & $7{,}80\cdot 10^{-4}$  & $-1{,}15\cdot 10^{4}$ \\
\bottomrule
\end{tabular}
\end{center}

\section{Elemente}
\label{sec:elem}

\subsection{Kantenelemente niedrigster Ordnung}

Verwendet werden Nédélec-Kantenelemente niedrigster Ordnung (Whitney-Formen).
Der Freiheitsgrad je Kante ist die Zirkulation
\begin{equation}
  a_e = \int_e \mathbf{A}\cdot\mathrm{d}\boldsymbol{\ell}
\end{equation}
entlang der gespeicherten Kantenrichtung. Der Raum enthält alle Felder der
Form $\mathbf{a} + \mathbf{b}\times\mathbf{r}$; ein homogenes
$\mathbf{B} = \nabla\times(\tfrac12\mathbf{B}\times\mathbf{r})$ ist damit
exakt darstellbar. Genau das ist der Patchtest in
Abschnitt~\ref{sec:verifikation}.

Unterstützt werden Tetraeder (6 Kanten), Hexaeder (12) und Prismen (9). Die
Element-zu-Kante-Tabelle hat einheitlich die Schrittweite 12; unbenutzte
Plätze tragen $-1$. Diese Konvention war die Quelle mehrerer Fehler
(Abschnitt~\ref{sec:fehler}).

\subsection{Orientierung}

Jede globale Kante wird als $(v_{\min}, v_{\max})$ abgelegt. Ein
Vorzeichen je Elementkante sagt, ob die lokale Richtung damit übereinstimmt.
Die lokalen Element\-matrizen sind \emph{unorientiert}; das Vorzeichen wird
erst beim Sammeln und Streuen angewandt:
\begin{equation}
  c_i^{\text{lokal}} = s_i\, a_{g(i)},
  \qquad
  b_{g(i)} \mathrel{+}= s_i\, F_i,
  \qquad
  A_{g(r),g(c)} \mathrel{+}= s_r s_c\, K_{rc}.
\end{equation}

\subsection{Abbildungen und Quadratur}

Der Code legt die Jacobi-Matrix \emph{transponiert} ab,
$J_{ij} = \partial x_j/\partial\xi_i$ (\texttt{compute\_jacobian} und die
Varianten für Tetraeder und Prisma), also $J = J_{\text{std}}^{\top}$ und
$J^{-1} = J_{\text{std}}^{-\top}$. Mit diesem $J$ gilt für die
Basisfunktionen die kovariante Piola-Transformation \emph{zeilenweise} über
$J^{-1}$, für ihre Rotation die kontravariante \emph{spaltenweise} über $J$
und geteilt durch $\det J$:
\begin{equation}
  \mathbf{W}^{\text{phys}} = J^{-1}\,\mathbf{W}^{\text{ref}}
  = J_{\text{std}}^{-\top}\,\mathbf{W}^{\text{ref}},
  \qquad
  \nabla\times\mathbf{W}^{\text{phys}}
  = \frac{1}{\det J}\,J^{\top}\left(\nabla\times\mathbf{W}^{\text{ref}}\right)
  = \frac{1}{\det J}\,J_{\text{std}}\left(\nabla\times\mathbf{W}^{\text{ref}}\right).
\end{equation}
Beide Formeln stehen nur einmal, in \texttt{piola\_covariant} und
\texttt{piola\_contravariant} (\texttt{NedelecDispatch.cuh}). Die Verwechslung
mit der transponierten Form ist \emph{nicht} schon auf affinen Elementen
folgenlos, sondern nur bei symmetrischem $J$, etwa auf einem achsenparallelen
Quader. Auf einem gedrehten und gescherten, aber affinen Element liefert die
transponierte Rotation aus den exakten Kantengraden eines homogenen Feldes
ein $\mathbf{B}$ mit $1{,}0$ relativem Fehler, die transponierte kovariante
Form ein $\mathbf{A}$ mit $0{,}66$ bis $0{,}68$ -- auf Hexaeder, Tetraeder
und Prisma gleichermaßen (\texttt{tests/test\_piola\_production},
Kontrollfall; $\|J - J^{\top}\|/\|J\| = 0{,}67$). Der Patchtest läuft deshalb
auf einem zufällig gestörten Netz, die Prüfung der Kraft- und Ausgabekörper
auf einem gedrehten und gescherten Element.

Quadratur: 4 Punkte auf dem Tetraeder, $2\times2\times2$-Gauss auf dem
Hexaeder. Der Tetraeder wird stets affin abgebildet, $\det J$ ist dort
konstant.

\section{Nullraum, Regularisierung, Projektionen}
\label{sec:regul}

Drei Eingriffe stehen zwischen der Assemblierung und dem Löser. Sie hängen
zusammen und werden hier in der Reihenfolge ihres Auftretens beschrieben.

\subsection{Regularisierung des Curl-Curl-Nullraums}

$K$ ist auf Gradientenfeldern exakt null ($\nabla\times\nabla\varphi = 0$).
In den nichtleitenden Gebieten -- Luft \emph{und} die stromgetriebenen Spulen,
die $\kappa = 0$ tragen -- ist $K + \mathrm{j}\omega C$ deshalb singulär.

Der Hauptpfad des Projekts löst das über eine Baum-Kobaum-Eichung mit
zusätzlichen Kohomologie-Projektoren (ein Volltorus-Loch je Ringspule). Dieser
Pfad wählt stattdessen eine Regularisierung:
\begin{equation}
  \sigma_{\text{Luft}} = \varepsilon\,\sigma_{\text{Leiter}},
  \qquad \varepsilon = 10^{-8} \text{ (Vorgabe, einstellbar)}.
\end{equation}
Das beseitigt beide Nullraumanteile auf einen Schlag. Der Preis ist eine
schlechtere Kondition und ein Artefakt, das Abschnitt~\ref{sec:schein}
beschreibt.

Die Regularisierung \emph{muss} vor der Assemblierung der Elementmatrizen
stehen, sonst ist sie wirkungslos, und die Wirksamkeit lässt sich am
Verhältnis $\max|\omega C| / \max|K|$ ablesen: ohne sie trüge nur die Platte
zu $C$ bei.

\subsection{Divergenzfreiheit des Quellterms}
\label{sec:quellterm}

$\Js$ wird punktweise vorgegeben ($J\hat{\boldsymbol{\varphi}}$ in den
Spulenelementen, null außerhalb) und dann integriert. Im Kontinuum ist
$\nabla\cdot\Js = 0$ und $\Js\cdot\mathbf{n} = 0$ auf der Spulenmantelfläche.
Diskret gilt das nicht: die Spulen werden \emph{elementweise} zugeordnet, ihre
Ersatzfläche ist gezackt, und deren Normalen stehen nicht senkrecht auf
$\hat{\boldsymbol{\varphi}}$. Der diskrete Spulenstrom ist damit nicht
geschlossen.

Gemessen mit dem diskreten Gradienten $G$ (Knoten $\to$ Kanten):
\begin{equation}
  \frac{\|G^{\!\top} b\|}{\|b\|} = 6{,}1\cdot 10^{-3},
  \qquad \text{Anforderung } \approx 10^{-4}.
\end{equation}

\textbf{Durch Netzverfeinerung ist das nicht behebbar.} Das ist gemessen, nicht
vermutet -- zweifach:

\begin{itemize}[nosep]
\item Im synthetischen Test (Abschnitt~\ref{sec:verifikation}) fällt der
  Fehler bei $\Js$ ohne Materialgrenze mit $h^{2{,}8}$, mit Materialgrenze nur
  mit $h^{0{,}38}$; auf einem unverzerrten Netz stagniert er praktisch.
\item Am echten Netz \emph{steigt} er unter \texttt{--refine=1} von
  $6{,}1\cdot10^{-3}$ auf $1{,}1\cdot10^{-2}$, obwohl die Geometrie
  unverändert bleibt (jede Facette wird in vier koplanare Dreiecke geteilt).
  Warum er steigt statt nur zu stagnieren, ist offen.
\end{itemize}

Behoben wird es durch orthogonale Projektion auf $\ker G^{\!\top}$:
\begin{equation}
  (G_F^{\!\top} G_F)\,\varphi = G_F^{\!\top} b,
  \qquad b' = b - G_F\varphi.
  \label{eq:proj}
\end{equation}
$G^{\!\top}G$ ist der Knoten-Graph-Laplace, gut konditioniert; rund 300
CG-Iterationen genügen. Zwei Details bestimmen die Umsetzung:

\begin{description}[style=nextline,leftmargin=1em]
\item[Konstanten.] $G^{\!\top}G$ ist singulär, konstantes $\varphi$ liegt im
  Kern. Die rechte Seite ist dazu automatisch orthogonal, das CG bleibt aber
  nur stabil, wenn der konstante Anteil \emph{während} der Iteration
  herausgenommen wird -- Rundung bringt ihn sonst mit Eigenwert 0 wieder ein.
\item[Freie Kanten.] Der Index $F$ in~\eqref{eq:proj} ist wesentlich. $b$ ist
  vor der Projektion nur im Spulenbereich besetzt, am Rand exakt null. $\varphi$
  löst aber ein \emph{globales} Poisson-Problem, also ist $G\varphi$ überall von
  null verschieden -- die Randkanten tragen danach einen Wert, den die
  Dirichlet-Behandlung zurücksetzt, und genau das bringt die Divergenz an den
  Randknoten wieder ein ($1{,}6\cdot10^{-5}$ statt $6\cdot10^{-16}$). Läuft die
  Projektion nur über die freien Kanten, bleibt $b'$ dort null per
  Konstruktion.
\end{description}

Ergebnis: $\|G^{\!\top}b\|/\|b\| = 6\cdot10^{-16}$ auf dem System, das
tatsächlich gelöst wird.

\subsection{Randbedingung}

$\mathbf{A}\times\mathbf{n} = 0$ auf dem äußeren Rand. Der Kantensatz gilt für
$n$ Kanten, im $2n$-System sind beide Hälften betroffen. Zeile und Spalte
werden genullt, die Diagonale auf einen Wert innerhalb der freien Diagonale
gesetzt.

\section{Löser}

\subsection{Direkt: cuDSS}
\label{sec:cudss}

Gelöst wird~\eqref{eq:block} als \texttt{CUDSS\_MTYPE\_GENERAL} mit
\texttt{MVIEW\_FULL}, Nested-Dissection-Ordnung und \texttt{PIVOT\_AUTO}. Die
symmetrische Variante wäre möglich, ist aber nicht umgesetzt.

Zwei Unverträglichkeiten von cuDSS~0.8.0 sind durch Bisektion über alle
Einzel- und Paarkombinationen gemessen, auf zwei Grafikkarten (sm\_120 und
sm\_75):
\begin{center}
\begin{tabular}{ll}
\toprule
Kombination & Ergebnis \\
\midrule
\texttt{DETERMINISTIC\_MODE} $+$ \texttt{IR\_N\_STEPS} $>0$ & ANALYSIS \texttt{NOT\_SUPPORTED} \\
\texttt{DETERMINISTIC\_MODE} $+$ \texttt{HYBRID\_MEMORY}    & ANALYSIS \texttt{NOT\_SUPPORTED} \\
\bottomrule
\end{tabular}
\end{center}
Ordnung und Pivotierung sind an beidem unbeteiligt; beide Male schränkt der
Determinismus die interne Ablaufplanung ein. Aufgelöst wird nach einer festen
Rangfolge: der Hybridmodus gewinnt immer, weil der Faktor sonst nicht in den
Gerätespeicher passt und gar nichts gerechnet wird; sonst gewinnt der
Determinismus über die Nachbesserung. Läufe mit Hybridmodus sind \emph{nicht}
bitweise reproduzierbar, und der Lauf sagt das ausdrücklich.

Größenordnungen auf dem feinsten Netz: $316\,476$ Zeilen, $1{,}03\cdot10^{7}$
Nichtnullen, Füllfaktor $43{,}3$, $4{,}46\cdot10^{8}$ Nichtnullen im Faktor.
Das sind 1,7\,GB allein für die Werte; auf einer 8-GB-Karte läuft es nur
mit Auslagerung.

\subsection{Iterativ: GMRES mit ILU(0)}
\label{sec:iterativ}

Zum Vergleich lässt sich dasselbe System iterativ lösen
(\texttt{--fem-direct-iterative}). GMRES bekommt dabei bitgenau dieselben
CSR-Arrays nach der Randbehandlung -- die einzige Variable ist der Löser. Ein
matrixfreier Weg über die Blockform wäre möglich, würde aber ein
\emph{anderes} System lösen, weil die Randbedingung dort an anderer Stelle
sitzt; ein Vergleich zweier Löser auf zwei Systemen sagt nichts.

Das Ergebnis ist zweigeteilt. \emph{Ohne} die Quellterm-Projektion stagniert
GMRES hart: von Iteration 400 bis 4000 bewegt sich das relative Residuum nicht
($2{,}58\cdot10^{-3}$). \emph{Mit} Projektion fällt es durch, aber langsam --
rund 1000 Iterationen je Dekade, $4{,}4\cdot10^{-6}$ nach 2000 Iterationen in
152\,s, nicht konvergiert.

Der Grund ist gemessen. Am Endresiduum $r$ gilt $\|G^{\!\top}r\|/\|r\| = 2{,}26$
bei $\|G^{\!\top}b\|/\|b\| = 6\cdot10^{-16}$ -- das Restresiduum ist reiner
Gradient, obwohl die rechte Seite dort nichts hat. Die Komponente kommt also
\emph{aus dem Operator}, nicht aus $b$: sie entsteht aus Rundung und wächst mit
den winzigen Eigenwerten von $\omega C$ auf den Gradientenfeldern -- Faktor 350
unter dem Rest, ablesbar an $\max|\omega C|/\max|K| = 2{,}85\cdot10^{-3}$.
ILU(0) hat dafür keinen Mechanismus.

Bemerkenswert und mit Vorsicht zu behandeln: die \emph{Kraft} ist trotz
Nichtkonvergenz identisch. Global weichen die Lösungen ab, in der Platte
stimmen sie auf unter ein Prozent überein, und $F_z$ ist in beiden Fällen
1,02\,N. Der stehende Anteil ist physikalisch träge, weil
$\nabla\times\nabla\varphi = 0$ und er die Platte nicht erreicht. Darauf
verlassen sollte man sich nicht -- ein Löser, dessen gemeldetes Residuum von
einem trägen Unterraum beherrscht wird, hat kein verwendbares
Abbruchkriterium. Die eigentliche Antwort ist ein Vorkonditionierer für den
Gradientenraum (AMS), und zwar auf dem inneren, reellen Operator statt auf der
$2n$-Blockform -- siehe Abschnitt~\ref{sec:ams-modus}.


\section{Warum reell und nicht komplex, und was das für AMS bedeutet}
\label{sec:ams}

Die Entscheidung gegen eine komplexe Systemgleichung ist keine Nebensache --
sie bestimmt, welche Vorkonditionierer überhaupt in Frage kommen. Dieser
Abschnitt nennt die Gründe, prüft sie gegen den naheliegenden Einwand und
korrigiert drei Punkte einer früheren Analyse.

\subsection{Die Gründe für die reelle Blockform}

\begin{enumerate}[nosep]
\item \textbf{Die vorhandene Maschinerie bleibt unverändert nutzbar}:
  \texttt{CoupledMVOperator}, ACA, ILU, AMS, \texttt{GMRESSolver}. Keine
  komplexen Kernel, keine komplexe ACA.
\item \textbf{Der Vorkonditionierer läuft über C-to-R} (Axelsson/Kucherov):
  für $H = \begin{bmatrix}W & -T\\ T & W\end{bmatrix}$ ist
  $P = \begin{bmatrix}W{+}T & 0\\ T{-}W & W{+}T\end{bmatrix}$, und eine
  Anwendung braucht zwei Lösungen mit $(W{+}T) = K + \omega C$. Das ist ein
  \emph{reeller} Operator von genau der Bauart, auf die AMS zielt. Kein neuer
  Vorkonditionierer, nur ein anderer Aufruf des vorhandenen.
\item \textbf{Gemessen} (\texttt{team28\_coarse\_12}, \texttt{--harmonic-study}):
  ohne C-to-R stagniert GMRES bei 4000 Iterationen, mit C-to-R sind es 38
  Iterationen auf $7{,}5\cdot10^{-9}$. Die Iterationszahl fällt um Faktor 105,
  während $\kappa_2$ sogar steigt -- nicht die Kondition entscheidet, sondern
  die Blockstruktur.
\end{enumerate}

\subsection{Korrektur zum iterativen Lauf dieses Berichts}

Der in Abschnitt~\ref{sec:iterativ} beschriebene Versuch wirft ILU(0) direkt
auf die rohe $2n$-Blockform. Genau das ist oben als stagnierend vermessen: die
Blockstruktur zahlt sich \emph{außen} aus, in C-to-R, nicht innen. Die dort
berichteten 4000 Iterationen bei $2{,}6\cdot10^{-3}$ sind damit keine neue
Erkenntnis, sondern die Reproduktion eines bekannten Ergebnisses. Der
begründete Weg wäre \texttt{solve\_harmonic} mit C-to-R gewesen; die
Rechtfertigung, dieselbe Matrix wie cuDSS zu verwenden, war formal richtig und
in der Abwägung falsch.

\subsection{Warum AMS auf dieser Form nicht ohne Weiteres anwendbar ist}

AMS (Hiptmair--Xu) glättet auf den Kanten und projiziert das verbleibende
Residuum über die diskrete Gradientenmatrix $G$ auf die Knoten, wo ein
Poisson-artiges Problem per AMG gelöst wird. Für die Blockform ergibt der
Knoten-Hilfsraum
\begin{equation}
  H_{\text{Knoten}}
  = \begin{bmatrix} G^{\!\top} & 0 \\ 0 & G^{\!\top}\end{bmatrix}
    \begin{bmatrix} K & -\omega C \\ \omega C & K \end{bmatrix}
    \begin{bmatrix} G & 0 \\ 0 & G \end{bmatrix}
  = \begin{bmatrix} 0 & -S \\ S & 0\end{bmatrix},
  \qquad S = \omega\,G^{\!\top}CG,
  \label{eq:hknoten}
\end{equation}
denn $KG = 0$ (die Rotation eines Gradienten verschwindet), also
$G^{\!\top}KG = 0$.

\paragraph{Die Nulldiagonale ist NICHT das Hindernis.} Eine frühere Analyse
liest~\eqref{eq:hknoten} als Sattelpunktstruktur, an der AMG scheitere. Das
trägt nicht:
\begin{equation}
  \begin{bmatrix} 0 & -S \\ S & 0\end{bmatrix}^{-1}
  \begin{bmatrix} a \\ b\end{bmatrix}
  = \begin{bmatrix} S^{-1}b \\ -S^{-1}a\end{bmatrix}.
\end{equation}
Die Nulldiagonale kostet nichts -- eine Anwendung sind zwei Lösungen mit $S$,
genau wie C-to-R außen zwei Lösungen mit $(W{+}T)$ braucht. Es ist kein
Sattelpunkt, sondern eine schiefsymmetrische Blockform mit invertierbarer
Nebendiagonale.

\paragraph{Das Hindernis ist, dass $S$ singulär ist.} $S$ entspricht
$\int \sigma\,\nabla\varphi\cdot\nabla\psi$, ist also ein $\sigma$-gewichteter
Poisson-Operator. Mit $\sigma = 0$ in der Luft hat $S$ einen Nullraum aus
allen $\varphi$, die dort leben -- derselbe Nullraum wie zuvor, nur von den
Kanten auf die Knoten verschoben.

Das ist richtig, und es ist trotzdem kein Ausschlusskriterium: der singuläre
Fall ist der \emph{Normalfall} des Verfahrens, nicht seine Grenze. Der Nullraum
wird nicht im skalaren Hilfsraum behandelt, sondern im vektoriellen.
Abschnitt~\ref{sec:ams-modus} führt das aus; der dort beschriebene Modus ist
inzwischen umgesetzt. Die Einschätzung in diesem Abschnitt war zu pessimistisch.

\paragraph{Und das überlebt den Wechsel zu komplex.} Für das komplexe System
ist der AMS-Knotenoperator
\begin{equation}
  G^{\!\top}(K + \mathrm{j}\omega C)\,G = \mathrm{j}\omega\,G^{\!\top}CG,
\end{equation}
weil $G^{\!\top}KG = 0$ unabhängig von der Arithmetik gilt. Dieselbe
Singularität, aus demselben Grund. Eine Standardbibliothek wie Hypre AMS
bräuchte dort ebenfalls eine Regularisierung oder eine Eichung. \textbf{Die
reelle Blockform ist also nicht die Ursache der AMS-Schwierigkeit}, und der
Wechsel zu komplex behebt sie nicht.

\subsection{Was der Wechsel zu komplex tatsächlich brächte}

Ein Argument bleibt, und es betrifft den direkten Löser. Die $2n$-Form
speichert jeden komplexen Eintrag als $2\times2$-Block, also vier reelle
Zahlen statt zwei, und zerlegt ein System doppelter Dimension:

\begin{center}
\begin{tabular}{lrrr}
\toprule
 & Zeilen & nnz & Faktor \\
\midrule
$2n$ reell & 316\,476 & $1{,}03\cdot10^{7}$ & $4{,}46\cdot10^{8}$ $\approx$ 3,6\,GB \\
komplex    & 158\,238 & $2{,}57\cdot10^{6}$ & grob $1{,}1\cdot10^{8}$ komplex $\approx$ 1,8\,GB \\
\bottomrule
\end{tabular}
\end{center}

Die zweite Zeile ist eine Schätzung mit gleichem Füllfaktor, keine Messung.
Der Faktor 2 ist aber genau das, was zwischen dem Projekt und dem
Determinismus steht: bei 1,8\,GB bräuchte es keinen Hybridspeicher, und
\texttt{DETERMINISTIC\_MODE} wäre wieder erlaubt (Abschnitt~\ref{sec:cudss}).
Ob cuDSS~0.8 \texttt{CUDA\_C\_64F} trägt, ist ungeprüft.

\subsection{Korrektur: die Regularisierung erzwingt keine Coulomb-Eichung}

Eine frühere Analyse begründet die Lösbarkeit damit, die Regularisierung
$\sigma_{\text{Luft}} = \varepsilon\sigma_{\text{Leiter}}$ forciere implizit
$\nabla\cdot\mathbf{A}\approx 0$. Das ist das Gegenteil des gemessenen
Verhaltens. Auf Gradientenfeldern lautet die Gleichung in der Luft
$\mathrm{j}\omega\varepsilon\sigma\,\nabla\varphi = -b_g$, also
\begin{equation}
  \nabla\varphi = -\frac{b_g}{\mathrm{j}\omega\,\varepsilon\,\sigma_{\text{Leiter}}}.
\end{equation}
Die Regularisierung macht das System lösbar, indem sie dem Gradientenanteil
einen um $1/\varepsilon$ \emph{verstärkten} Wert zuweist, nicht indem sie ihn
unterdrückt. Gemessen wächst max$|\Ai|$ von $0{,}0671$ auf $671$, wenn
$\varepsilon$ von $10^{-6}$ auf $10^{-10}$ fällt -- exakt linear in
$1/\varepsilon$ über vier Dekaden (Abschnitt~\ref{sec:schein}).

\subsection{Umgesetzt: AMS ohne Regularisierung (\texttt{--fem-direct-ams})}
\label{sec:ams-modus}

Der Abschnitt oben schließt, das Hindernis sei die Singularität von
$S = \omega\,G^{\!\top}CG$. Das stimmt, ist aber kein Hindernis, sondern der
\emph{Normalfall} des Verfahrens. Hypre kennt dafür den singulären Modus
(\texttt{HYPRE\_AMSSetBetaPoissonMatrix(NULL)}): der skalare Hilfsraum wird
dort stillgelegt, wo er nichts trägt, und die Gradientenmoden übernimmt der
\emph{vektorielle} Hilfsraum über den Alpha-Poisson-Operator mit
$\nu = 1/\mu_0 > 0$ überall. Die Ausblendung ist im Projekt seit längerem
vorhanden (\texttt{ams\_mask\_meaningless\_rows}); was fehlte, war die
Entscheidung, sie auch so zu benutzen.

Damit lässt sich die Regularisierung ersatzlos streichen. Der Modus
\texttt{--fem-direct-nullspace=ams} tut das:

\begin{itemize}[nosep]
\item Kein $\sigma_{\text{Luft}}$, keine Baum-Kobaum-Eichung. Das System bleibt
  singulär und konsistent.
\item Die $2n$-Matrix wird nicht gebaut. \texttt{solve\_harmonic} wendet $H$
  matrixfrei über $W = K + \omega C$ und $C$ an; C-to-R braucht je Anwendung
  zwei Lösungen mit $W$, und das ist der reelle Operator, auf den AMS zielt.
\item Der Direktlöser ist ausgeschlossen und der iterative Weg erzwungen. Auf
  einer echt singulären Matrix pivotiert cuDSS ins Nichts, ohne dass
  \texttt{CUDSS\_DATA\_INFO} es meldet -- auf diesem Pfad ist eine singuläre
  Matrix schon einmal mit \texttt{info = 0} zerlegt worden.
\item Die Quellterm-Projektion ist Pflicht statt Option. Konsistenz heißt
  $G^{\!\top}b = 0$, und Abschnitt~\ref{sec:quellterm} liefert dafür
  $6\cdot10^{-16}$.
\end{itemize}

Zwei Punkte der Umsetzung sind keine Nebensache.

\paragraph{Der Alpha-Poisson-Operator musste erst auf Tetraeder gebracht
werden.} Er war hexaederfest -- acht Knoten, trilineare Formfunktionen -- und
las auf einem Tetraedernetz Knotenplätze, auf denen $-1$ steht. Ohne ihn ist
der vektorielle Hilfsraum auf den Netzen dieses Pfades gar nicht aufbaubar.
Gemessener AMG-Konvergenzfaktor je V-Zyklus auf strukturierten Netzen mit
Leiter/Luft-Grenze:

\begin{center}
\begin{tabular}{lrrr}
\toprule
 & $8^3$ & $12^3$ & $16^3$ \\
\midrule
Hexaeder  & 0,248 & 0,264 & 0,278 \\
Tetraeder & 0,300 & 0,375 & 0,373 \\
\bottomrule
\end{tabular}
\end{center}

Zum Vergleich liegt der 7-Punkt-Laplace bei 0,30. Der Tetraederfall ist
zwischen den beiden feinsten Netzen flach; der Anstieg davor ist der Vorsprung
des groben Netzes, dessen Hierarchie nur drei Ebenen tief ist und dessen
Grobgitterproblem fast direkt gelöst wird.

\paragraph{Die Dirichlet-Bedingung braucht in $W$ und $C$ verschiedene
Diagonalwerte.} Festgehalten wird im $n$-Raum, in beiden Matrizen, und zwar mit
\begin{equation}
  W_{ee} = d, \qquad (\omega C)_{ee} = 0 .
\end{equation}
Nur dann lautet die festgehaltene Zeile $[d, 0]$ beziehungsweise $[0, d]$, und
Real- und Imaginärteil der Randkante werden unabhängig voneinander auf null
gesetzt. Bekäme $C$ dort denselben Wert $d$, stünde da $[d - \omega d,\,
-\omega d]$: eine Gleichung, die den Imaginärteil an den Realteil \emph{koppelt}.
Das System bliebe lösbar, es wäre nur ein anderes, und sichtbar würde das
allein als seltsame Iterationszahl.
\texttt{tests/test\_harmonic\_dirichlet\_split} prüft es an der Form, die der
Löser wirklich anwendet, und enthält die Gegenprobe mit dem absichtlich
eingebauten Fehler.

\paragraph{Gemessen.} Netz „Platte fein", 133\,899 Tetraeder,
158\,238 Kanten, RTX 3090.

\begin{center}
\begin{tabular}{lrr}
\toprule
 & $\varepsilon$ + cuDSS & AMS ohne Regularisierung \\
\midrule
$\max|B_r|$ Platte      & 0,0719\,T   & 0,072\,T \\
$F_z$                   & 1,020\,N    & 1,0\,N \\
Beitrag der Nichtleiter & $-3{,}83\cdot10^{-5}$\,N & $-2{,}2\cdot10^{-15}$\,N \\
\bottomrule
\end{tabular}
\end{center}

Die Abnahme war ausdrücklich \emph{nicht} das Residuum -- die
Baum-Kobaum-Eichung aus Abschnitt~\ref{sec:baum} erreichte
$4{,}9\cdot10^{-14}$ und lieferte trotzdem ein anderes Feld. Geprüft ist $B$,
und $B$ stimmt.

Der Löser braucht 15 äußere Iterationen auf ein relatives Residuum von
$5\cdot10^{-9}$, dafür 30 innere Lösungen mit $(W+T)$ -- genau zwei je äußerer
Iteration, wie C-to-R es verlangt -- mit zusammen 3995 Iterationen, im Mittel
133. Laufzeit 1100\,s.

Zwei Zahlen sind wichtiger als diese:

\begin{itemize}[nosep]
\item Das Residuum wurde \emph{unabhängig} nachgerechnet, im physikalischen
  Raum aus $W$ und $C$ statt im Tilde-Raum matrixfrei:
  $1{,}4\cdot10^{-8}$ gegen $5\cdot10^{-9}$. Die beiden Wege entstehen
  getrennt; ihre Übereinstimmung schließt einen Fehler in der Äquilibrierung
  oder in der Blockanwendung aus.
\item Der Beitrag der Nichtleiter ist $-2{,}2\cdot10^{-15}$\,N. Ohne
  Regularisierung ist $\sigma$ dort exakt null, also $J = 0$ und die Differenz
  reine Rundung. Die Nichtleiter-Maske ist in diesem Modus nachweislich
  redundant.
\end{itemize}

\paragraph{Ein Fehler auf dem Weg dorthin, weil er lehrreich ist.} Der erste
Lauf blieb ab Iteration 60 bei einem relativen Residuum von 0,058 stehen --
völlig flach über 150 weitere Iterationen. Ursache war die Richtung der
Äquilibrierung: die rechte Seite wurde durch $D$ geteilt statt multipliziert.

Das ist keine bloße Skalierung. Der Operator im Tilde-Raum ist
$\tilde A = DAD$, sein Nullraum also $D^{-1}n$, und
\begin{equation}
  \underbrace{(Db)\cdot(D^{-1}n) = b\cdot n = 0}_{\text{richtig}},
  \qquad
  \underbrace{(b/D)\cdot(D^{-1}n) = b^{\!\top}D^{-2}n \neq 0}_{\text{falsch}} .
\end{equation}
Die Quelltermprojektion hatte $G^{\!\top}b$ im physikalischen Raum auf
$5{,}79\cdot10^{-16}$ gebracht; die falsche Richtung trug den Gradientenanteil
danach wieder hinein, gewichtet mit $D^{-2}$. Im regularisierten Modus fällt
das nicht auf -- dort ist $A$ regulär und jede rechte Seite konsistent.

Das völlig flache Plateau ist die Signatur: ein Krylov-Verfahren kann den
Anteil außerhalb des Bildraums prinzipiell nicht verkleinern. Der Lauf prüft
die Konsistenz jetzt im Tilde-Raum, bevor er anfängt.

\paragraph{Noch offen: die Netzreihe.} Ein Messpunkt zeigt, dass das Verfahren
richtig rechnet, nicht dass es skaliert. Ob die äußere Iterationszahl über die
Netzreihe ungefähr konstant bleibt, ist die eigentliche Zusage von AMS und
steht noch aus.

\subsection{Der verbleibende Weg: projiziertes GMRES}

Bleibt die Deflation: den Gradientenanteil nicht nur aus $b$, sondern aus
\emph{jedem} Krylow-Vektor entfernen, über
$v' = v - G\varphi$ mit $(G^{\!\top}G)\varphi = G^{\!\top}v$.

Dieser Weg ist ausdrücklich \emph{nicht} durch die Messung in
Abschnitt~\ref{sec:iterativ} widerlegt. Dort wurde gezeigt, dass die
Projektion von $b$ die Konvergenz nicht verbessert -- weil die
Gradientenkomponente aus dem Operator kommt und nicht aus der rechten Seite.
Jeden Krylow-Vektor zu projizieren greift genau dort an: es entfernt den
Nullraum aus dem Suchraum statt nur aus der Quelle. Der Unterschied zwischen
einer einmaligen und einer fortlaufenden Projektion ist hier der ganze Punkt.

Zwei Vorbehalte. Ein Projektor ist singulär und lässt sich nicht als
gewöhnlicher Rechts-Vorkonditionierer einhängen; es braucht ein projiziertes
Verfahren, GMRES auf $PHP\,y = Pb$ im Bildraum von $P$ -- ein Eingriff in den
Löser, nicht ein zusätzlicher Aufruf. Und es fehlt die Kantenglättung des
vollen AMS, die schlechte Kondition des Curl-Curl-Operators im
rotationsbehafteten Raum bleibt also. Erwartbar ist damit das Ende der
Stagnation, nicht eine kleine Iterationszahl.

\subsection{Verworfen: die Baum-Kobaum-Eichung}
\label{sec:baum}

Der naheliegende Ausweg -- statt zu regularisieren wieder zu eichen -- wurde
versucht und verworfen. Die Topologie erlaubt es: mit mitvernetzter Luft ist
das nichtleitende Gebiet Luft plus Spulen, also Kasten minus Platte, und damit
einfach zusammenhängend. Im BEM-Fall war es die Spule selbst, ein Volltorus,
daher dort ein Kohomologie-Generator je Spule.

Umgesetzt (\texttt{--fem-direct-gauge}, 16\,521 von 158\,238 Kanten geeicht)
ergab sich ein Residuum von $4{,}9\cdot10^{-14}$ -- das System ist also
regulär, die Eichung greift. Aber:
\begin{center}
\begin{tabular}{lrr}
\toprule
 & regularisiert & geeicht \\
\midrule
Residuum & $4{,}3\cdot10^{-15}$ & $4{,}9\cdot10^{-14}$ \\
global max$|\Ar|$ & $6{,}65\cdot10^{-5}$ & $4{,}24\cdot10^{-4}$ \\
$B_r$ Platte & 0,0719\,T & 0,0736\,T \\
$F_z$ & 1,02\,N & 1,25\,N \\
\bottomrule
\end{tabular}
\end{center}

$B$ ist eichinvariant -- eine Abweichung von 2,4\,\% bedeutet also nicht einen
anderen Eichrepräsentanten, sondern ein anderes Feld. Als Ursache bleibt die
Kondition: ein kleines Residuum ist bei schlecht konditionierter Matrix kein
kleiner Fehler, und der Baum wird über Union-Find bei aufsteigenden
Kantenindizes gebaut, was lange Pfade quer durchs Netz ergibt. Eine
Breitensuche von der Leiterkomponente aus wäre die bekannte Gegenmaßnahme,
ist aber nicht umgesetzt -- der Weg wurde zugunsten von AMS aufgegeben.

Nebenbefund derselben Messung: die Hälfte der Quelltermnorm
($1{,}32\cdot10^{4}$ von $2{,}61\cdot10^{4}$) liegt auf Baumkanten und wird
verworfen, weil der Baum das nichtleitende Gebiet spannt und dazu die
stromgetriebenen Spulen gehören. Zulässig ist das für ein konsistentes $b$
(die Baumzeilen sind dann redundant), aber es ist unschön und war der erste
Verdacht für die Abweichung.

\section{Kraftauswertung}

\subsection{Volumenintegral (Lorentz)}

Zeitgemittelt und in Phasoren:
\begin{equation}
  \langle \mathbf{J}\times\mathbf{B}\rangle
  = \frac{\sigma\omega}{2}\left(
      \Ai\times(\nabla\times\Ar) - \Ar\times(\nabla\times\Ai)
    \right),
\end{equation}
integriert über die Leiter. Eine Besonderheit: der Ausschluss der Nichtleiter
kann \emph{nicht} über $\kappa > 0$ erfolgen, weil die Regularisierung jedem
Element $\kappa > 0$ gibt. Stattdessen wird vor der Regularisierung eine Maske
der betroffenen Elemente angelegt. Der Beitrag des regularisierten
Luftgebietes ist gemessen $1{,}2\cdot10^{-7}$~N gegen ein Ergebnis von
1,02\,N -- er wäre also auch ohne Maske vernachlässigbar, ist aber
$\varepsilon$-unabhängig und deshalb von der $\varepsilon$-Reihe nicht
aufdeckbar.

\subsection{Oberflächenintegral (Maxwell) als Gegenprobe}

Als unabhängiger Weg zur selben Größe der Spannungstensor über die
Materialgrenzflächen. Zeitgemittelt:
\begin{equation}
  \langle \mathsf{T}\cdot\mathbf{n}\rangle
  = \frac{1}{2\mu_0}\left[
      (\mathbf{B}_r\!\cdot\!\mathbf{n})\mathbf{B}_r
    + (\mathbf{B}_i\!\cdot\!\mathbf{n})\mathbf{B}_i
    - \tfrac12\left(|\mathbf{B}_r|^2 + |\mathbf{B}_i|^2\right)\mathbf{n}
    \right].
\end{equation}
Die Normale wird geometrisch bestimmt (Flächenmitte gegen Elementmitte), nicht
aus einer Tabelle -- die Umlaufsinnkonvention ist je Elementtyp anders und ein
Fehler dort dreht nur das Vorzeichen.

\textbf{Dieser Weg trägt auf den vorliegenden Netzen nicht}, und das ist
gemessen. Mit $\Ai = 0$ verschwindet der Realteil der Stromdichte in der
Platte; ohne Strom innerhalb der Hülle und bei $\mu_r = 1$ muss das Integral
über eine geschlossene Fläche exakt null ergeben. Es tut es nicht:

\begin{center}
\begin{tabular}{rrrr}
\toprule
Grenzflächen & Nullprobe [N] & relativ zu $F_z$ & Ordnung \\
\midrule
  208 & $-4{,}06$  & $4{,}82$ & --- \\
  832 & $-1{,}18$  & $0{,}94$ & $h^{1{,}78}$ \\
13\,274 & $-0{,}492$ & $0{,}48$ & $h^{0{,}63}$ \\
\bottomrule
\end{tabular}
\end{center}

Der Kernel selbst ist in Ordnung -- bei homogenem $\mathbf{B}$ liefert er auf
einer geschlossenen Hülle $1{,}09\cdot10^{-16}$ relativ. (Nachtrag
31.10.2026: auf \emph{Viereck}flächen -- Hexaeder, Seitenflächen des
Prismas -- fehlte bis dahin der Faktor 4 der Flächenabbildung; die Hülle der
Prüfung besteht nur aus Dreiecken, und die Zahlen dieses Abschnitts stammen
von reinen Tetraedernetzen, sind also nicht betroffen.
\texttt{tests/test\_piola\_production} prüft seither je Fläche.) Die Aufgabe ist ihm
zu schwer, aus zwei Gründen: die Nettokraft entsteht als Differenz von
Einzelbeiträgen um $|B|^2/(2\mu_0)$ mal Hüllfläche, hier rund 200\,N;
und das diskrete $\mathbf{B}$ ist über Elementflächen \emph{nicht} stetig, weil
Nédélec nur die Tangentialstetigkeit von $\mathbf{A}$ und damit die
Normalstetigkeit von $\mathbf{B}$ zusichert. Der Spannungstensor braucht beide
Komponenten. Bei einer 3\,mm dünnen Platte aus flachen Tetraedern ist der
Sprung besonders groß.

\section{Verifikation}
\label{sec:verifikation}

Vier Testprogramme laufen ohne GPU (\texttt{ctest -L host}).

\begin{center}
\begin{tabular}{llr}
\toprule
Prüfung & Programm & Messwert \\
\midrule
Konstantes $\mathbf{A}$ exakt reproduziert & \texttt{test\_patch\_assembly} & $4{,}4\cdot10^{-16}$ \\
Patchtest $\nabla\times\mathbf{A}_h = \mathbf{B}$ & \texttt{test\_patch\_assembly} & $2{,}3\cdot10^{-15}$ \\
Konstantes $\Js$: $\|G^{\!\top}b\|/\|b\|$ innen & \texttt{test\_patch\_assembly} & $1{,}25\cdot10^{-16}$ \\
Maxwell auf geschlossener Hülle & \texttt{test\_patch\_assembly} & $1{,}09\cdot10^{-16}$ \\
Verfeinerung: Volumen erhalten & \texttt{test\_mesh\_refinement} & $1{,}2\cdot10^{-14}$ \\
Projektion: reiner Gradient entfernt & \texttt{test\_source\_projection} & $5{,}2\cdot10^{-12}$ \\
Projektion: Idempotenz & \texttt{test\_source\_projection} & $2{,}4\cdot10^{-14}$ \\
\bottomrule
\end{tabular}
\end{center}

Zwei Punkte verdienen Hervorhebung. Der Patchtest läuft \emph{über
Elementgrenzen}, auf einem Netz aus Kuhn-Tetraedern mit zufällig verschobenen
inneren Knoten, und sammelt die Koeffizienten aus den globalen
Freiheitsgraden. Damit sind erstmals die globale Kantennummerierung, die
Orientierungsvorzeichen und die Konformität mitgeprüft -- ein Vorzeichenfehler
in dieser Schicht erzeugt keinen Absturz, sondern ein falsches Feld, und ein
Einzelelementtest sieht ihn nicht.

Die dritte Zeile ist die schärfste: bei konstantem $\Js$ ist der Integrand auf
einem affinen Tetraeder linear und wird von der 4-Punkt-Regel exakt
integriert. Das Ergebnis muss also \emph{maschinengenau} null sein, nicht nur
klein. Damit ist ein Vorzeichen- oder Nummerierungsfehler als Ursache der
$6{,}1\cdot10^{-3}$ aus Abschnitt~\ref{sec:quellterm} ausgeschlossen.

\section{Messergebnisse}

\subsection{Netze}

\begin{center}
\begin{tabular}{lrrrrl}
\toprule
Netz & Knoten & Tetraeder & Kanten & Platte & Plattenlage \\
\midrule
Startlage        & 766    & 3\,692   & 4\,709   & 259     & $z = 3{,}8\ldots6{,}8$ mm \\
Gleichgewicht    & 1\,759 & 8\,216   & 10\,683  & 259     & $z = 11{,}3\ldots14{,}3$ mm \\
\quad$+$\texttt{--refine=1} & 12\,442 & 65\,728 & 81\,005 & 2\,072 & dito \\
Platte fein      & 23\,249 & 133\,899 & 158\,238 & 19\,812 & dito \\
\bottomrule
\end{tabular}
\end{center}

Die Spulen sind in allen Netzen dieselben 517 und 251 Elemente. Daraus folgte
eine überprüfbare Vorhersage, die eingetroffen ist: $\|G^{\!\top}b\|/\|b\|$,
$\|b\|$ und die Zahl der besetzten Einträge sind auf allen Netzen identisch.
Die Verschlechterung unter \texttt{--refine=1} ist damit eindeutig der
Verfeinerung \emph{der Spulen} zuzuschreiben.

\subsection{Levitationskraft}

\begin{center}
\begin{tabular}{lrrr}
\toprule
Netz & Platte & $F_z$ [N] & gegen Gewicht 1,026\,N \\
\midrule
Startlage ($z_{\text{Mitte}} = 5{,}3$ mm) & 259 & $1{,}97$ & --- \\
Gleichgewicht & 259 & $0{,}842$ & $-17{,}9\,\%$ \\
$+$\texttt{--refine=1} & 2\,072 & $1{,}250$ & $+21{,}8\,\%$ \\
Platte fein & 19\,812 & $\mathbf{1{,}020}$ & $\mathbf{-0{,}6\,\%}$ \\
\bottomrule
\end{tabular}
\end{center}

Das Plattengewicht ist aus Geometrie, Dichte und $g$ gerechnet, nichts daran
ist angepasst. Der Sollwert aus der Projektdokumentation lautet 1,0497\,N
bei 11,3\,mm; der Startlagenwert von 1,97\,N liegt darüber, wie es
sein muss -- die Platte ist dort 7,5\,mm näher an den Spulen.

\textbf{Die Übereinstimmung von $0{,}6\,\%$ ist noch nicht belegt.} Die Reihe
schwingt, und zwischen den beiden letzten Netzen änderte sich nur die
Auflösung in der Ebene (beide zwei Lagen), was $F_z$ um $-18{,}4\,\%$
verschob. Ein fünfter Punkt mit feinerer Ebene würde entscheiden, ob
1,02\,N Konvergenz oder Zufall ist.

Bemerkenswert dabei: das \emph{Feld} ist konvergiert, die Kraft war es nicht.
Zwischen den beiden letzten Netzen bleibt $\max|B_r|$ in der Platte auf
$1{,}6\,\%$ stehen, während $F_z$ um $48\,\%$ sprang. Unteraufgelöst war also
das Kraftintegral, nicht das Feld -- 259 Tetraeder in einer 3\,mm dünnen
Scheibe sind eine einzige Lage, und der Wirbelstrom hat über die Dicke ein
Profil, das eine Lage nicht darstellen kann.

\subsection{Unempfindlichkeit gegen \texorpdfstring{$\varepsilon$}{eps}}
\label{sec:schein}

Die Regularisierung erzeugt ein Artefakt, das lange die Genauigkeit des
gesamten Pfades bestimmte. Auf Gradientenfeldern ist $K$ exakt null, die
Blockform~\eqref{eq:block} reduziert sich dort auf $x_r = 0$,
$x_i = -b_g/(\omega C)$: der Gradientenanteil von $b$ landet vollständig im
Imaginärteil, verstärkt um $1/\varepsilon$.

\begin{center}
\begin{tabular}{lrrrr}
\toprule
$\varepsilon$ & \multicolumn{2}{c}{ohne Projektion} & \multicolumn{2}{c}{mit Projektion} \\
 & Residuum & $\max|\Ai|$ & Residuum & $\max|\Ai|$ \\
\midrule
$10^{-6}$  & $1{,}26\cdot10^{-12}$ & $0{,}0671$ & $4{,}36\cdot10^{-15}$ & $2{,}38\cdot10^{-6}$ \\
$10^{-7}$  & $1{,}20\cdot10^{-11}$ & $0{,}671$  & ---                   & --- \\
$10^{-8}$  & $1{,}25\cdot10^{-10}$ & $6{,}71$   & $4{,}31\cdot10^{-15}$ & $2{,}38\cdot10^{-6}$ \\
$10^{-9}$  & $1{,}26\cdot10^{-9}$  & $67{,}1$   & ---                   & --- \\
$10^{-10}$ & $1{,}21\cdot10^{-8}$  & $671$      & $4{,}33\cdot10^{-15}$ & $2{,}38\cdot10^{-6}$ \\
\bottomrule
\end{tabular}
\end{center}

Ohne Projektion skaliert $\max|\Ai|$ exakt mit $1/\varepsilon$ über vier
Dekaden, und das Residuum wächst mit -- letzteres aus dem Dynamikumfang: eine
Lösung, die von $6{,}65\cdot10^{-5}$ bis $6{,}88$ spannt, lässt sich in
doppelter Genauigkeit nicht auf $10^{-15}$ auflösen.

Mit Projektion bewegt sich keine Ziffer. Das Artefakt ist um Faktor
$2{,}9\cdot10^{6}$ kleiner, das Residuum um fünf Dekaden besser, und
$\varepsilon$ ist über vier Dekaden ohne Wirkung -- auf \emph{jede} Größe des
Laufs, nicht nur auf $\mathbf{B}$.

Wichtig für die Einordnung: die physikalischen Größen waren schon \emph{vor}
der Projektion unempfindlich. $\max|\Ar|$ und $\max|\Ai|$ in der Platte,
$\mathbf{B}$ und $F_z$ blieben über die gesamte $\varepsilon$-Reihe
unverändert. Das Artefakt saß vollständig in der Luft, wo
$\nabla\times\nabla\varphi = 0$ es unschädlich macht. Die Projektion beseitigt
also einen Fehler, der die Kraft nie erreicht hat -- ihr Nutzen liegt in der
Konditionierung und in der Erfüllung der Divergenzanforderung, nicht in einer
korrigierten Kraft.

\section{Fehler, die keinen Absturz erzeugten}
\label{sec:fehler}

Ein eigener Abschnitt, weil das Muster übertragbar ist. Sieben Fehler lagen
zwischen dem ersten Aufruf und der ersten Kraftzahl, und keiner davon führte
zu einem Absturz oder einer Fehlermeldung.

\begin{enumerate}[nosep]
\item \textbf{Uninitialisierte Stackplätze.} Bei einem Tetraeder bleiben die
  Kantenplätze 6--11 ungeschrieben, die Multiplikation lief aber über alle
  zwölf. Die Matrix ist dort sauber null -- und $0{,}0 \cdot \mathrm{NaN}$ ist
  NaN. 90 NaN in der rechten Seite, laufabhängig, weil nur ein Teil der
  Threads ein NaN-Bitmuster im Frame erwischt.
\item \textbf{Drei Indexzugriffe auf $-1$.} Das Nullsetzen des
  \emph{Vorzeichens} neutralisiert den Wert, nicht den Index.
\item \textbf{Unerreichbare Regularisierung.} Die Vorbedingungsprüfung
  beanstandete $\kappa = 0$ als harten Fehler und lief vor dem Block, der
  genau das beseitigt.
\item \textbf{Wirkungslose Regularisierung.} Danach erreichbar, aber die
  Assemblierung berechnete die Elementmatrizen nicht neu -- die
  Materialänderung erreichte die Matrix nicht. Folge war eine singuläre
  Matrix, die cuDSS mit \texttt{CUDSS\_DATA\_INFO = 0} zerlegte.
\item \textbf{Rechte Seite als Zeitprobe.} Faktor $\sin(\omega\Delta t) =
  0{,}0314$ statt der Amplitude; jede Kraft hätte am gewählten Zeitschritt
  gehangen.
\item \textbf{Entwerteter Nichtleiter-Test.} Nach der Regularisierung hat
  jedes Element $\kappa > 0$, die Abfrage im Kraftintegral trennte nichts
  mehr.
\item \textbf{Falsche Entwarnung in der Diagnose.} Der Ausdruck
  \texttt{(b2 > 0.0) ? sqrt(g2/b2) : 0.0} fällt bei
  \texttt{b2 = NaN} in den else-Zweig -- jeder Vergleich mit NaN ist falsch --
  und meldete daraufhin ausdrücklich ``divergenzfrei, Inkonsistenz
  ausgeschlossen'' für eine vollständig NaN-verseuchte rechte Seite.
\end{enumerate}

Der letzte Punkt ist der lehrreichste: eine Diagnose, die genau dann Entwarnung
gibt, wenn etwas kaputt ist, ist schlechter als keine.

\section{Offene Punkte}

\begin{description}[style=nextline,leftmargin=1em]
\item[Abschneidefaktor.] Der Luftkasten reicht radial nur bis Faktor $1{,}745$
  der Leiterausdehnung, verlangt sind mindestens 2. Die Randbedingung
  $\mathbf{A}\times\mathbf{n} = 0$ verformt das Feld damit dort, wo integriert
  wird. Dieser systematische Anteil steckt in allen vier Kraftwerten und ist
  in keinem quantifiziert -- auch nicht im Vorzeichen. Er braucht ein Netz mit
  radial größerem Kasten.
\item[Fünfter Netzpunkt.] Zwei Lagen, Ebene feiner, alles andere unverändert.
  Er entscheidet, ob die $0{,}6\,\%$ Konvergenz oder Zufall sind.
\item[Vorkonditionierung.] Siehe Abschnitt~\ref{sec:ams}. Standard-AMS ist
  hier nicht anwendbar, und zwar nicht wegen der reellen Blockform, sondern
  weil der AMS-Knotenoperator $\omega G^{\!\top}CG$ mit $\sigma = 0$ in der
  Luft singulär ist -- was der Wechsel zu einer komplexen Formulierung nicht
  behebt. Es bleiben zwei Wege: C-to-R mit ILU im inneren
  $(K+\omega C)$-Löser, gemessen 38 Iterationen auf einem gröberen Netz, aber
  auf diesem Netz nicht erprobt; und projiziertes GMRES (Deflation), das die
  Stagnation beenden dürfte, ohne die Kondition im rotationsbehafteten Raum zu
  verbessern. Die Baum-Kobaum-Eichung ist versucht und verworfen
  (Abschnitt~\ref{sec:baum}).
\item[Zweiter Kraftpfad.] Der Maxwell-Spannungstensor konvergiert, aber aus zu
  großer Entfernung; zwei weitere Verfeinerungsstufen brächten ihn auf rund
  $9\,\%$. Auf der vorhandenen Hardware ist das nicht zu haben. Solange
  bleiben die 1,02\,N ohne unabhängige Bestätigung.
\item[Determinismus.] Läufe mit Hybridspeicher sind nicht bitweise
  reproduzierbar. Für das feinste Netz ist der Hybridmodus nicht verhandelbar,
  also gilt \texttt{test\_determinism\_fingerprint} für diesen Pfad nicht.
\end{description}

\end{document}
